\renewcommand{\theposttitle}{How to overcome your impossible standards}
\renewcommand{\thepostdate}{2024-05-29}
\vspace*{\fill}
\begin{center}
\includegraphics[width=\textwidth, height=0.55\textheight, keepaspectratio]{C:/Users/James/Documents/GitHub/VisualIdeasBook/data/images/503f6b35b967b6ae47a9c1f99a5367e3.png}
\end{center}
\vspace{2mm}
\begin{center}
\begin{minipage}{0.92\textwidth}
\small
\textbf{Avoid the second mistake.}\\
Inspired by James Clear’s book, \textbf{Atomic Habits}.

Now the trick is how?

The best piece of advice I’ve heard from James Clear is, “reduce the scope to maintain the habit.”

Most of my life I’ve held almost impossible to meet standards, for myself and others. While this was useful to me for a season, I realized over time those standards distorted my creative output into an all or nothing approach.

The harsh truth is that \textbf{a balanced life doesn’t mean balanced days.} Balance comes in the form of an averaging of your days.

There will be days when I only get a few push-ups or a brisk walk in for exercise. There are days when I get home so late that I only have enough time with the kids for a short cuddle and bedtime story.

The point is to embrace the imperfect messiness of it all. Miss a day, you’re human. Try your best to get back on the bike day 2, even if it means you don’t go as far forward as you’d like.

A little progress is still progress.
\end{minipage}
\end{center}
\vspace*{\fill}